% title:          Reflection property of the ellipse
% area:           differential-geometry, physical-proofs
% methods:
% difficulty:
% difficulty-ai:  easy
% status:
% review:         none
% --- history: all optional, leave EMPTY if unknown ---
% origin:         historical
% origin-date:
% origin-number:  Book III, Proposition 48
% also-in:        book
% taken-from:
% references:     Apollonius of Perga, Conics, Book III, Proposition 48 (late 3rd century BC): the focal distances of a point make equal angles with the tangent - T. L. Heath, Apollonius of Perga: Treatise on Conic Sections (Cambridge, 1896), Proposition 71 [III. 48], p. 116 - https://archive.org/download/treatiseonconics00apolrich/treatiseonconics00apolrich_djvu.txt; Wikipedia, Ellipse, section 'Focus-to-focus reflection property' - https://en.wikipedia.org/wiki/Ellipse; taken from M. Levi, The Mathematical Mechanic: Using Physical Reasoning to Solve Problems (Princeton UP, 2009), Section 3.1 'The Optical Property of Ellipses', p. 29 (statement word for word; mechanical proof with a weighted pulley on a string, p. 30) - https://books.google.ch/books?id=2Jp3FKRcZbEC&jscmd=SearchWithinVolume2&q=Let%20P%20be%20a%20point%20on%20an%20ellipse&scoring=r; second proof with a ring sliding on the ellipse attached to the foci by two constant-tension springs, Section 3.2 'More about the Optical Property', p. 31 - https://books.google.ch/books?id=2Jp3FKRcZbEC&jscmd=SearchWithinVolume2&q=ring%20springs&scoring=r; legacy set note: physics-based problem (force balance / potential energy)
% history:        ai
\begin{problem}
Let $P$ be a point on an ellipse with the foci $F_1$ and $F_2$,
and let $MN$ be the tangent at $P$. Prove that 
\[ \angle F_1PM=\angle F_2PN\]
\end{problem}
